%*###############################################################
\section{Numerical Simulation Results of EIT Cooling}
\label{sec:eit_cooling_results}
%*###############################################################

The numerical model described in Section~\ref{sec:numerical_eit} was evaluated to characterize the Electromagnetically Induced Transparency (EIT) cooling dynamics of $^{87}\mathrm{Rb}$ atoms confined in the optical conveyor belt. To systematically distinguish the fundamental cooling mechanism from the effects of the complex atomic structure, the simulations were performed in two stages. First, an idealized three-level $\Lambda$ system was simulated to validate the analytical predictions. Subsequently, a complete 24-level model incorporating the full hyperfine and Zeeman structure of the $D_2$ line was implemented to assess realistic experimental conditions, including off-resonant excitation and population leakage.

%*###############################################################
\subsection{Ideal three-level system dynamics}
\label{subsec:results_three_level}
%*###############################################################

The steady-state spectroscopic response of the ideal three-level system was obtained by scanning the probe-laser detuning $\Delta_p$ while keeping the coupling field fixed. The resulting total excited-state population exhibits the characteristic Fano profile associated with EIT, as shown in Figure~\ref{fig:fano_full_3} and detailed in Figure~\ref{fig:fano_3}.

\begin{figure}[htbp]
\centering
\includegraphics[width=0.8\textwidth]{Cap_3_figures/fano_full_3.png}
\caption{Steady-state excited population profile of the idealized three-level system as a function of the probe detuning $\Delta_p$. The transparency window (dashed red line) corresponds to the two-photon resonance condition, leading to the formation of a dark state.}
\label{fig:fano_full_3}
\end{figure}

The coupling laser intensity was set according to the exact cooling condition (Eq.~\eqref{eq:exact_condition_numerical}), tuning the AC Stark shift so that the narrow absorption maximum exactly aligns with the red motional sideband (cooling transition), while the transparency minimum coincides with the carrier transition.

\begin{figure}[htbp]
\centering
\includegraphics[width=0.8\textwidth]{Cap_3_figures/fano_3.png}
\caption{Detailed EIT profile of the three-level system, showing the alignment of the laser detunings with the motional sidebands. The narrow absorption maximum is precisely aligned with the red sideband (solid red line) to maximize the cooling rate, while the carrier (dashed green line) falls in the transparency minimum, suppressing heating.}
\label{fig:fano_3}
\end{figure}

The time evolution of the average phonon number $\langle n \rangle$ for this configuration is presented in Figure~\ref{fig:time_evol_3}. Starting from an initial state with a high mean phonon number, the system undergoes rapid ground-state cooling. The average phonon number decays exponentially, approaching the theoretical minimum limit near $\langle n \rangle \approx 0$ in less than $0.1\,\mathrm{ms}$. This confirms that, in the absence of competing decay channels, the targeted sideband asymmetry drives highly efficient motional cooling.

\begin{figure}[htbp]
\centering
\includegraphics[width=0.8\textwidth]{Cap_3_figures/time_evol_3.png}
\caption{Time evolution of the average phonon number $\langle n \rangle$ in the three-level simulation. The population rapidly cools to the motional ground state within approximately $0.1\,\mathrm{ms}$.}
\label{fig:time_evol_3}
\end{figure}

%*###############################################################
\subsection{Multilevel simulation and optical pumping effects}
\label{subsec:results_multilevel}
%*###############################################################

When extending the simulation to the full 24-level structure of the $^{87}\mathrm{Rb}$ $D_2$ line, the dynamics become significantly more complex due to spontaneous emission into ground states outside the designed $\Lambda$ subsystem.

To demonstrate this effect, the steady-state spectrum was first simulated without an active repump field (Figure~\ref{fig:fano_24_norep}). In this scenario, the excited-state population drops to near zero across the entire frequency scan. Following a few excitation cycles, the atoms decay into uncoupled Zeeman sublevels (such as those in the $F=2$ or $F=1$ manifolds not addressed by the probe and coupling beams). This optical pumping completely depletes the active EIT levels, effectively turning off the cooling mechanism and resulting in the loss of the atomic ensemble.

\begin{figure}[htbp]
\centering
\includegraphics[width=0.8\textwidth]{Cap_3_figures/fano_24_norep.png}
\caption{Steady-state excited population of the 24-level system without a repump field. Optical pumping causes the population to accumulate in dark, uncoupled Zeeman states, causing the spectroscopic signal to vanish.}
\label{fig:fano_24_norep}
\end{figure}

To prevent this population leakage, a $\pi$-polarized repump field was added to recycle the atoms back into the EIT cycle. The steady-state profile obtained with the optimized repump field is shown in Figure~\ref{fig:fano_24}. While the Fano profile and the dark resonance are recovered, the spectrum deviates from the ideal three-level case. The absorption peak is broader, and the transparency minimum at the carrier frequency is less deep. These imperfections are a direct consequence of off-resonant coupling to adjacent excited manifolds (such as $F'=3$) and the continuous perturbation caused by the repump field, which introduces additional scattering events that weakly break the ideal dark state.

\begin{figure}[htbp]
\centering
\includegraphics[width=0.8\textwidth]{Cap_3_figures/fano_24.png}
\caption{Excited state population profile of the 24-level system with an active repump field. The Fano profile is restored, allowing alignment of the red sideband with the absorption peak. The resonance is broader and the carrier suppression is less ideal compared to the three-level model, due to off-resonant scattering and repump dynamics.}
\label{fig:fano_24}
\end{figure}

The impact of this complex internal structure on the cooling dynamics was evaluated using Monte Carlo trajectory simulations (\texttt{mcsolve}). Figure~\ref{fig:time_evol_24} shows the time evolution of the average phonon number for the 24-level system. EIT sideband cooling is still achieved, ultimately bringing the atoms toward the motional ground state. However, the effective cooling rate is drastically reduced compared to the ideal case.

\begin{figure}[htbp]
\centering
\includegraphics[width=0.8\textwidth]{Cap_3_figures/time_evol_24.png}
\caption{Time evolution of the average phonon number $\langle n \rangle$ in the complete 24-level system (average of 150 quantum trajectories). Cooling to $\langle n \rangle \ll 1$ is achieved, but the required time increases to nearly $18\,\mathrm{ms}$ due to population recycling and off-resonant scattering.}
\label{fig:time_evol_24}
\end{figure}

The system takes nearly $18\,\mathrm{ms}$ to approach $\langle n \rangle \approx 0$, which is roughly two orders of magnitude longer than the three-level simulation. This slowdown occurs because the atoms spend a significant fraction of their time outside the optimal $\Lambda$ system, waiting to be recycled by the repump laser. Furthermore, the excess scattering from the repumper and the imperfect carrier suppression introduce recoil heating events that must be continuously compensated by the sideband cooling mechanism. Despite this reduction in cooling speed, the simulated cooling timescale is fully compatible with state-of-the-art EIT cooling experiments implemented in free-space, where similar multilevel effects are present~\cite{Chow2024}.